%%%%%%%%%%  scan idx 408  |  folio 401  %%%%%%%%%%

S G A  7


\underline{E X P O S E  XXII}


\underline{\underline{UNE FORMULE DE CONGRUENCE POUR LA FONCTION}} $\zeta$

\underline{par N. KATZ}


0. \underline{Introduction}

Soit $X$ un schéma propre sur $\mathbb F_q$. Sa fonction zêta, considérée
comme élément de $\mathbb Z[[T]]$, peut-être réduite modulo $p$ ; on veut la
calculer modulo $p$. On trouve une expression du type auquel on pouvait
s'attendre depuis les formules de points fixes du type ``Woods Hole'', comme
produit alterné de polynômes caractéristiques d'un ``Frobenius'' convenable
agissant sur l'un quelconque des trois espaces de cohomologie suivants :

1) la cohomologie cohérente $H^i(X,\mathcal O_X)$ (3.1.1)

2) la cohomologie de De Rham
$H^i_{\mathrm{DR}}(X)=\mathbb H^i(X,\Omega_X^\bullet)$ (3.1.1)

3) la cohomologie étale
$H^i_{\mathrm{et}}(X,\mathbb Z/p\mathbb Z)$ (3.1.2).

Cette expression doit être vue comme une amélioration d'un cas spécial de la
formule des traces de Woods Hole : celui du faisceau structurel et de
l'endomorphisme de Frobenius. Ce cas est, en effet, un corollaire de la
formule de congruence (cf. 3.2.1).

La méthode de cet exposé est complètement naïve ; elle consiste à se ramener
au cas d'une hypersurface, ensuite à ``réduire modulo $p$''

%%%%%%%%%%  scan idx 409  |  folio 402  %%%%%%%%%%

le calcul de DWORK pour la vraie fonction zêta d'une hypersurface. En
``réduisant'' soigneusement, on trouve l'amélioration 5.0.6 due à AX [1]
(et également à DWORK dans le cas lisse [4, \S 7])du théorème bien connu de
WARNING-CHEVALLEY [8].

Une autre méthode, entièrement cohomologique et dans l'esprit de SGA 5,
s'appliquant également à des ``coefficient tordus'', vient d'être trouvée par
DELIGNE, et devrait figurer dans la réédition de SGA 5 ; elle repose sur la
formule des traces type ``Woods Hole'', et la``formule de Künneth symétrique''
de DELIGNE (SGA 4 XVII 5.5.21).

1. \underline{Un rappel sur les opérations} ``$p$-\underline{linéaires}''

1.0. Soit $V$ un espace vectoriel de dimension finie sur un corps $k$ de
caractéristique $p>0$. Désignons par $q$ une puissance de $p$. Una application

\begin{equation}\tag{1.0.1}
\varphi:V\longrightarrow V
\end{equation}

s'appelle $q$-\underline{linéaire} si elle additive et si elle vérifie

\begin{equation}\tag{1.0.2}
\varphi(\lambda v)=\lambda^q\varphi(v)
\qquad \forall\lambda\in k,\ v\in V.
\end{equation}

1.0.3. On fera attention au fait que l'itéré $n$'ième $\varphi^n$ d'un
endomorphisme $q$-linéaire est $q^n$-linéaire, et que, si $k$ n'est pas
parfait, l'image $\varphi(V)$ d'un endomorphisme $q$-linéaire n'est pas
nécessairement un $k$-sous-espace vectoriel.

1.0.4. Pour chaque extension $K/k$, un endomorphisme $q$-linéaire $\varphi$
de $V$ donne lieu à un unique endomorphisme $q$-linéaire $\varphi_K$ de
$V_K=V\otimes_k K$, satisfaisant à la formule

\begin{equation}\tag{1.0.5}
\varphi_K(\mu\otimes v)=\mu^q\varphi(v)
\qquad\text{pour }\mu\in K,\ v\in V.
\end{equation}

%%%%%%%%%%  scan idx 410  |  folio 403  %%%%%%%%%%

1.0.6. On dit qu'un endomorphisme $q$-linéaire $\varphi$ de $V$ est
semi-simple si son image \underline{engendre} $V$ comme $k$-vectoriel, et
qu'il est nilpotent s'il admet un itéré nul.

On définit des $k$-sous-espace stables par $\varphi$

\begin{equation}\tag{1.0.7}
V_{\mathrm{ss}}=
\bigcap_{n\geq1}
\bigl(\text{le $k$-sous-espace engendré par }\varphi^n(V)\bigr)
\end{equation}

\begin{equation}\tag{1.0.8}
V_{\mathrm{nilp}}=
\bigcup_{n\geq1}\bigl(\text{le noyau de }\varphi^n:V\longrightarrow V\bigr).
\end{equation}

Évidemment, $\varphi$ restreint à $V_{\mathrm{ss}}$ est semi-simple, et
$\varphi$ restreint à $V_{\mathrm{nilp}}$ est nilpotent. \underline{Si} $k$ \underline{est parfait},
on a

\begin{equation}\tag{1.0.9}
V=V_{\mathrm{ss}}\oplus V_{\mathrm{nilp}},
\end{equation}

mais, si $k$ n'est pas parfait, (1.0.9) devient inexact, car il y a des
endomorphismes $q$-linéaires qui sont injectifs sans être semi-simples.

Pour une extension $K/k$ quelconque, on a

\begin{equation}\tag{1.0.10}
(V_K)_{\mathrm{ss}}=V_{\mathrm{ss}}\otimes_k K.
\end{equation}

\underline{Proposition} 1.1. \underline{Soit $\varphi$ un endomorphisme
$q$-linéaire d'un espace vectoriel $V$ de dimension finie sur un corps
séparablement clos $k$. Si $\varphi$ est semi-simple alors, posant}

\[
V^{1-\varphi}=\text{\underline{les $\mathbb F_q$-vectoriel formé des points fixes de}
$\varphi$}=\Ker(1-\varphi),
\]

\underline{on a}

\begin{equation}\tag{1.1.1}
V\xleftarrow{\ \sim\ }
(V^{1-\varphi})\otimes_{\mathbb F_q}k.
\end{equation}

%%%%%%%%%%  scan idx 411  |  folio 404  %%%%%%%%%%

\underline{Corollaire} 1.1.2.
\[
V^{1-\varphi^n}\xleftarrow{\ \sim\ }
(V^{1-\varphi})\otimes_{\mathbb F_q}\mathbb F_{q^n}.
\]

\underline{Démonstration.} (LANG, cf. [6], p. 119).

Soit $\underline{e}$ une $k$-base de $V$, et écrivons

\begin{equation}\tag{1.1.3}
\varphi(\underline{e})=A\underline{e},\qquad A\in\mathrm{GL}(n,k)
\end{equation}

Désignons par

\begin{equation}\tag{1.1.4}
F:\mathrm{GL}(n,k)\longrightarrow\mathrm{GL}(n,k)
\end{equation}

la flèche ``élévation des coordonnées à la $q$'ième puissance''. Alors pour
$g\in\mathrm{GL}(n,k)$, $g\underline{e}$ est une base de $V$ formée de points fixes de
$\varphi$ si et seulement si

\begin{equation}\tag{1.1.5}
F(g)Ag^{-1}=1.
\end{equation}

Désignons par $\bar k$ une clôture algébrique de $k$. Le
\underline{groupe} $\mathrm{GL}(n,\bar k)$ \underline{agit} sur la variété
$\mathrm{GL}(n,\bar k)$ par

\begin{equation}\tag{1.1.6}
\left\{
\begin{aligned}
\mathrm{GL}(n,\bar k)\times\mathrm{GL}(n,\bar k)
&\longrightarrow\mathrm{GL}(n,\bar k)\\
(g,A)&\longmapsto F(g)Ag^{-1}.
\end{aligned}
\right.
\end{equation}

Pour $A$ fixé, la flèche

\begin{equation}\tag{1.1.7}
g\longmapsto F(g)Ag^{-1}
\end{equation}

est étale comme on voit en calculant son application tangente, donc les
orbites sont ouvertes.Comme $\mathrm{GL}(n,\bar k)$ est irréductible, il n'y
a qu'une orbite ; autrement dit, la flèche

\begin{equation}\tag{1.1.8}
\begin{aligned}
\mathrm{Gl}(n,\bar k)&\longrightarrow\mathrm{Gl}(n,\bar k)\\
g&\longrightarrow F(g^{-1})\cdot g
\end{aligned}
\end{equation}

%%%%%%%%%%  scan idx 412  |  folio 405  %%%%%%%%%%

est surjective. Parce qu'elle est aussi étale, il s'ensuit, $k$ étant
séparablement clos que la flèche

\begin{equation}\tag{1.1.9}
\begin{aligned}
\mathrm{GL}(n,k)&\longrightarrow\mathrm{GL}(n,k)\\
g&\longrightarrow F(g^{-1})\cdot g
\end{aligned}
\end{equation}

est surjective, ce qui veut dire que $V$ admet une $k$-base formée des points
fixes de $\varphi$. Il est alors clair, $\varphi$ étant $q$-linéaire, que tout
point fixe de $\varphi$ se trouve dans le $\mathbb F_q$-vectoriel engendré par
une $k$-base formée des points fixes de $\varphi$, ce qui prouve 1.1.

\underline{Corollaire} 1.1.10. \underline{Soit $\varphi$ un endomorphisme
$q$-linéaire d'un espace vectoriel $V$ de dimension finie sur un corps
séparablement clos $k$. Alors on a un isomorphisme canonique}

\begin{equation}\tag{1.1.11}
V_{\mathrm{ss}}\xleftarrow{\ \sim\ }
V^{1-\varphi}\otimes_{\mathbb F_q}k.
\end{equation}

Il suffit en effet d'appliquer 1.1 à $V_{\mathrm{ss}}$.

\underline{Proposition} 1.2. \underline{Soit $\varphi$ un endomorphisme
$q$-linéaire d'un espace vectoriel de dimension finie sur un corps
séparablement clos $k$. Alors l'application additive}

\[
1-\varphi:V\longrightarrow V
\]

\underline{est surjective.}

\underline{Démonstration.} On a une suite exacte de groupes abéliens

\[
0\longrightarrow V_{\mathrm{ss}}\longrightarrow V
\longrightarrow V/V_{\mathrm{ss}}\longrightarrow0,
\]

%%%%%%%%%%  scan idx 413  |  folio 406  %%%%%%%%%%

sur laquelle agit $1-\varphi$. Par le ``lemme du serpent'', il suffit de
traiter $V_{\mathrm{ss}}$, évident par 1.1 en employant une base formée des
points fixes, et $V/V_{\mathrm{ss}}$, sur lequel $\varphi$ est nilpotente,
donc sur lequel $1-\varphi$ est inversible.

2. \underline{Cohomologie cohérente et cohomologie étale}

\underline{Proposition} 2.0. \underline{Soit} $\bar X$ \underline{un schéma propre sur
un corps} $k$ \underline{séparablement clos, et désignons par}

\begin{equation}\tag{2.0.1}
\overline{\mathfrak F}_p:H^{\cdot}(\bar X,\mathcal O_{\bar X})
\longrightarrow H^{\cdot}(\bar X,\mathcal O_{\bar X})
\end{equation}

\underline{l'application} $p$-\underline{linéaire induite par}

\begin{equation}\tag{2.0.2}
\overline{\mathfrak F}_p:\mathcal O_{\bar X}\longrightarrow\mathcal O_{\bar X},
\qquad \overline{\mathfrak F}_p(f)=f^p.
\end{equation}

\underline{Désignons par}

\[
H^{\cdot}_{\mathrm{et}}(\bar X,\mathbb Z/p\mathbb Z)
\]

\underline{les groupes de cohomologie étales à coefficients dans le faisceau
constant} $\mathbb Z/p\mathbb Z$. \underline{Alors on a des isomorphismes canoniques :}

\begin{equation}\tag{2.0.3}
H^i_{\mathrm{et}}(\bar X,\mathbb Z/p\mathbb Z)
\simeq
(H^i(\bar X,\mathcal O_{\bar X}))^{1-\overline{\mathfrak F}_p}.
\end{equation}

\underline{Démonstration.} On a une suite exacte de faisceaux sur
$\bar X_{\mathrm{et}}$.

\begin{equation}\tag{2.0.4}
0\longrightarrow\mathbb Z/p\mathbb Z\longrightarrow\mathcal O_{\bar X}
\xrightarrow{\ 1-\mathfrak F_p\ }\mathcal O_{\bar X}\longrightarrow0
\end{equation}

qui donne une suite exacte de cohomologie

\begin{equation}\tag{2.0.5}
\cdots\longrightarrow H^i_{\mathrm{et}}(\bar X,\mathbb Z/p\mathbb Z)
\longrightarrow H^i_{\mathrm{et}}(\bar X,\mathcal O_{\bar X})
\xrightarrow{\ 1-\mathfrak F_p\ }
H^i_{\mathrm{et}}(\bar X,\mathcal O_{\bar X})\longrightarrow\cdots .
\end{equation}

On a
$H^i_{\mathrm{et}}(\bar X,\mathcal O_{\bar X})\simeq
H^i(\bar X,\mathcal O_{\bar X})$ ; c'est un $k$-vectoriel \underline{de dimension finie}
$\bar X$ étant propre. En appliquant 1.2, on trouve des suites exactes

%%%%%%%%%%  scan idx 414  |  folio 407  %%%%%%%%%%

\begin{equation}\tag{2.0.6}
0\longrightarrow H^i_{\mathrm{et}}(\bar X,\mathbb Z/p\mathbb Z)
\longrightarrow H^i(\bar X,\mathcal O_{\bar X})
\xrightarrow{\ 1-\mathfrak F_p\ }
H^i(\bar X,\mathcal O_{\bar X})\longrightarrow0,
\end{equation}

cqfd. Utilisant 1.1.10, on conclut de (2.0) :

\underline{Corollaire} 2.1. \underline{Pour} $\bar X$ \underline{propre sur un corps
séparablement clos} $k$,

\begin{equation}\tag{2.1.1}
H^i(\bar X,\mathcal O_{\bar X})_{\mathrm{ss}}
\xleftarrow{\ \sim\ }
H^i_{\mathrm{et}}(\bar X,\mathbb Z/p\mathbb Z)
\otimes_{\mathbb F_p}k.
\end{equation}

2.2. Soit maintenant $f:X\longrightarrow\Spec(\mathbb F_q)$ un schéma
propre sur $\mathbb F_q$. Désignons par $k$ une clôture algébrique de
$\mathbb F_q$, et posons

\begin{equation}\tag{2.2.1}
\bar X=X\times_{\mathbb F_q}k.
\end{equation}

Le faisceau étale $R^if_*(\mathbb Z/p\mathbb Z)$ sur
$\operatorname{Sp}(\mathbb F_q)$ s'identifie à
$H^i(\bar X,\mathbb Z/p\mathbb Z)$, vu comme
$\Gal(k/\mathbb F_q)$ module. Désignons par

\begin{equation}\tag{2.2.2}
\mathfrak F_q:H^i(X,\mathcal O_X)\longrightarrow H^i(X,\mathcal O_X)
\end{equation}

l'opération linéaire (elle est $q$-linéaire, mais on est sur $\mathbb F_q$!)
induite par

\begin{equation}\tag{2.2.3}
\mathfrak F_q:\mathcal O_X\longrightarrow\mathcal O_X,
\qquad \mathfrak F_q(f)=f^q.
\end{equation}

\underline{Proposition} 2.2.4. \underline{Désignons par}
$\varphi_q\in\Gal(k/\mathbb F_q)$ \underline{l'élément de} ``\underline{Frobenius}'',
$\varphi_q(\lambda)=\lambda^q$ \underline{et soit} $i\in\mathbb Z$. \underline{On a une identité de
polynômes caractéristiques}

\begin{equation}\tag{2.2.4.1}
\det_{\mathbb F_p}\!\left(
1-t\varphi_q^{-1}\mid H^i_{\mathrm{et}}(\bar X,\mathbb Z/p\mathbb Z)
\right)
=
\det_{\mathbb F_q}\!\left(
1-t\mathfrak F_q\mid H^i(X,\mathcal O_X)
\right).
\end{equation}

\underline{où le premier membre est le polynôme caractéristique de l'élément}
$\varphi_q^{-1}$ \underline{dans} $\pi_1(\operatorname{Sp}(\mathbb F_q))$ \underline{agissant sur le
module galoisien} $H^i(\bar X,\mathbb Z/p\mathbb Z)$.

En effet, 2.2.4 est un corollaire de la

%%%%%%%%%%  scan idx 415  |  folio 408  %%%%%%%%%%

\underline{Proposition} 2.2.5. \underline{Soit} $X$ \underline{propre sur}
$\Spec(\mathbb F_q)$. \underline{Il y a un isomorphisme canonique}

\begin{equation}\tag{2.2.6}
H^i(X,\mathcal O_X)_{\mathrm{ss}}\otimes_{\mathbb F_q}k
\simeq
H^i_{\mathrm{et}}(\bar X,\mathbb Z/p\mathbb Z)
\otimes_{\mathbb F_p}k
\end{equation}

\underline{qui transforme l'automorphisme}
$\mathfrak F_q\otimes_{\mathbb F_q}k$ \underline{du premier membre en l'automorphisme}
$\varphi_q^{-1}\otimes_{\mathbb F_q}k$ \underline{du second} (\underline{qu'on peut aussi interpréter
comme l'endomorphisme} $\mathrm{fr}_q^{X*}\otimes_{\mathbb F_p}k$, \underline{où}
$\mathrm{fr}_q^X:\bar X\longrightarrow\bar X$ \underline{est le} $k$-\underline{endomorphisme de
Frobenius.} cf. SGA 5, XVIII, 2.3.4).

\underline{Démonstration.} Compte tenu de ce que

\begin{equation}\tag{2.2.7}
H^i(\bar X,\mathcal O_{\bar X})\simeq
H^i(X,\mathcal O_X)\otimes_{\mathbb F_q}k,
\end{equation}

on a, par 1.0.10,

\begin{equation}\tag{2.2.8}
H^i(\bar X,\mathcal O_{\bar X})_{\mathrm{ss}}\simeq
H^i(X,\mathcal O_X)_{\mathrm{ss}}\otimes_{\mathbb F_q}k,
\end{equation}

de sorte qu'on obtient l'isomorphisme 2.2.6 en composant les isomorphismes
2.2.8 et 2.1.1.

Compte tenu de 2.2.7 l'opération $q$-linéaire

\begin{equation}\tag{2.2.9}
\mathfrak F_q:H^i(X,\mathcal O_X)\longrightarrow H^i(X,\mathcal O_X)
\end{equation}

admet une extension $q$-linéaire (cf. 1.0.5))

\begin{equation}\tag{2.2.10}
\overline{\mathfrak F}_q:H^i(\bar X,\mathcal O_{\bar X})
\longrightarrow H^i(\bar X,\mathcal O_{\bar X}),
\end{equation}

et admet aussi une extension $k$-linéaire, notée encore
$\mathfrak F_q$

\begin{equation}\tag{2.2.11}
\mathfrak F_q:H^i(\bar X,\mathcal O_{\bar X})
\longrightarrow H^i(\bar X,\mathcal O_{\bar X}).
\end{equation}

%%%%%%%%%%  scan idx 416  |  folio 409  %%%%%%%%%%

On désigne par

\begin{equation}\tag{2.2.12}
\varphi_q:H^i(\bar X,\mathcal O_{\bar X})
\longrightarrow H^i(\bar X,\mathcal O_{\bar X})
\end{equation}

l'opération $q$-linéaire dont les pointes fixes sont $H^i(X,\mathcal O_X)$.
Les opérateurs $\mathfrak F_q$, $\varphi_q$ et
$\overline{\mathfrak F}_q$, moyennant (2.2.7), s'écrivent aussi
$\mathfrak F_q\otimes\mathrm{id}_k$, $\mathrm{id}\otimes\varphi_q$,
$\mathfrak F_k\otimes\varphi_k$, donc commutent entre eux, et on a

\begin{equation}\tag{2.2.13}
\overline{\mathfrak F}_q
=\varphi_q\mathfrak F_q
=\mathfrak F_q\circ\varphi_q,
\end{equation}

et ils commutent tous à $\overline{\mathfrak F}_p$, de sorte que $\mathfrak F_q$ et
$\varphi_q$ agissent sur
$H^i(\bar X,\mathcal O_{\bar X})^{1-\overline{\mathfrak F}_q}$ et sur
$H^i(\bar X,\mathcal O_{\bar X})^{1-\overline{\mathfrak F}_p}$. Or,
$\mathfrak F_q$ et $\varphi_q$ sont des inverses l'une à l'autre \underline{sur}
$H^i(\bar X,\mathcal O_{\bar X})^{1-\overline{\mathfrak F}_q}$, par (2.2.13).
Donc l'opération de $\mathfrak F_q$ sur

\[
H^i(\bar X,\mathcal O_{\bar X})_{\mathrm{ss}}
\simeq
H^i(\bar X,\mathcal O_{\bar X})^{1-\overline{\mathfrak F}_q}
\otimes_{\mathbb F_q}k
\qquad\text{(cf. 1.1.11)}
\]

est l'inverse de l'opération de $\varphi_q\otimes k$ sur

\[
H^i(\bar X,\mathcal O_{\bar X})^{1-\overline{\mathfrak F}_q}
\otimes_{\mathbb F_q}k.
\]

Écrivant

\begin{equation}\tag{2.2.14}
H^i(\bar X,\mathcal O_{\bar X})^{1-\mathfrak F_q}
\simeq
H^i(\bar X,\mathcal O_{\bar X})^{1-\mathfrak F_p}
\otimes_{\mathbb F_p}\mathbb F_q,
\end{equation}

l'automorphisme $\varphi_q$ de
$H^i(\bar X,\mathcal O_{\bar X})^{1-\mathfrak F_q}$ devient, via
2.2.14, l'automorphisme $\varphi_q\otimes\mathbb F_q$ de
$H^i(\bar X,\mathcal O_{\bar X})^{1-\overline{\mathfrak F}_p}
\otimes_{\mathbb F_p}\mathbb F_q$, de sorte que l'automorphisme
$\varphi_q\otimes k$ de $H^i(\bar X,\mathcal O_{\bar X})_{\mathrm{ss}}$
devient, via 2.1.1, l'automorphisme $\varphi_q\otimes k$ de
$H^i_{\mathrm{et}}(\bar X,\mathbb Z/p\mathbb Z)\otimes_{\mathbb F_p}k$,
ce qui achève la démonstration (2.2.5).

%%%%%%%%%%  scan idx 417  |  folio 410  %%%%%%%%%%

3. \underline{La formule de congruence : énoncés et équivalences élémentaires}

Je tiens à remercier Grothendieck et Deligne pour m'avoir expliqué les
variantes (6.1.2), (3.1.1.2) et (3.1.3) de (3.1.1),

3.0. Soit $X$ un schéma de type fini sur $\mathbb F_q$. On définit le degré
d'un point fermé $y\in X$ comme le degré de son corps résiduel
$\mathbb F_q(y)/\mathbb F_q$. On définit la fonction zêta de
$X/\mathbb F_q$ comme élément de $\mathbb Z[[t]]$, par la formule
(cf. aussi SGA 5 XIII 3.1)

\begin{equation}\tag{3.0.1}
Z(t,X/\mathbb F_q)=\prod_y(1-t^{\deg(y)})^{-1}.
\end{equation}

Par la ``fonction zêta modulo $p$'' on veut dire toujours l'image de
$Z(t,X/\mathbb F_q)$ dans $\mathbb F_q[[t]]$ par la flèche de ``réduction
modulo $p\mathbb Z[[t]]$''

\[
\mathbb Z[[t]]\longrightarrow\mathbb F_q[[t]].
\]

\underline{Théorème} 3.1. \underline{Soit} $X$ \underline{un schéma propre sur}
$\mathbb F_q$. \underline{Alors}

\begin{equation}\tag{3.1.1}
Z(t,X/\mathbb F_q)
\equiv
\prod_{i\geq0}
\det\!\left(1-t\mathfrak F_q\mid H^i(X,\mathcal O_X)\right)^{(-1)^{i+1}}
\pmod p
\end{equation}

\underline{ou, ce qui est équivalent par} (2.2.4).

\begin{equation}\tag{3.1.2}
Z(t,X/\mathbb F_q)
\equiv
\prod_{i\geq0}
\det\!\left(1-t\varphi_q^{-1}\mid
H^i_{\mathrm{et}}(\bar X,\mathbb Z/p\mathbb Z)\right)^{(-1)^{i+1}}
\pmod p
\end{equation}

\underline{Corollaire} 3.2. \underline{Soit} $X$ \underline{un schéma propre sur}
$\mathbb F_q$.

\underline{Alors}

\begin{equation}\tag{3.2.1}
\sum_{i\geq0}(-1)^i
\operatorname{trace}\!\left(\mathfrak F_q\mid H^i(X,\mathcal O_X)\right.
\equiv \#X(\mathbb F_q)\pmod p
\end{equation}

\underline{ou, ce qui est équivalent par} (2.2.4)

%%%%%%%%%%  scan idx 418  |  folio 411  %%%%%%%%%%

\begin{equation}\tag{3.2.2}
\sum_{i\geq0}(-1)^i
\operatorname{tr}\!\left(\varphi_q^{-1}\mid
H^i_{\mathrm{et}}(\bar X,\mathbb Z/p\mathbb Z)\right)
\equiv \#X(\mathbb F_q)\mod p.
\end{equation}

On prend les coefficients de $t$ dans les congruences (3.1.1) et (3.1.2).

\underline{Théorème} 3.3. \underline{Soit} $X$ \underline{un schéma propre sur}
$\mathbb F_q$ \underline{on désigne par} $H^i_{\mathrm{DR}}(X/\mathbb F_q)$ \underline{les}
$\mathbb F_q$-\underline{vectoriels} $\mathbb H^i(X,\Omega^\bullet_{X/\mathbb F_q})$, \underline{sur lesquels}
$\mathfrak F_q$ \underline{agit}. \underline{On a}

\begin{equation}\tag{3.3.1}
Z(t,X/\mathbb F_q)
\equiv
\prod_{i\geq0}
\det\!\left(1-t\mathfrak F_q\mid
H^i_{\mathrm{DR}}(X/\mathbb F_q)\right)^{(-1)^{i+1}}
\mod p.
\end{equation}

En effet, utilisant la suite spectrale

\begin{equation}\tag{3.3.2}
E_1^{p,q}=H^q(X,\Omega^p_{X/\mathbb F_q})
\Longrightarrow H^{p+q}_{\mathrm{DR}}(X/\mathbb F_q)
\end{equation}

sur laquelle $\mathfrak F_q$ agit (via l'endomorphisme
$\mathfrak F_q=\mathrm{fr}_p^{X*}$ déduit de
$\mathrm{fr}_q^X:X\longrightarrow X$) en tuant $E_1^{p,q}$ pour $q\ne0$, on
voit que (3.3.1) équivaut à (3.1.1).

Donnons un petit complément.

La formation de la ``partie semi-simple'' (cf. 1.0.7) étant un foncteur
exact, on déduit de (3.3.2) une suite spectrale

\begin{equation}\tag{3.3.3}
(E_1^{p,q})_{\mathrm{ss}}
\Longrightarrow H^{p+q}_{\mathrm{DR}}(X/\mathbb F_q)_{\mathrm{ss}}.
\end{equation}

Comme $(E_1^{p,q})_{\mathrm{ss}}=0$ pour $q\ne0$, on en déduit pour tout
$i$ un isomorphisme

\begin{equation}\tag{3.3.4}
H^i_{\mathrm{DR}}(X/\mathbb F_q)_{\mathrm{ss}}
\xrightarrow{\ \sim\ }
H^i(X,\mathcal O_X)_{\mathrm{ss}},
\end{equation}

commutant à $\mathfrak F_q$, d'où

\begin{equation}\tag{3.3.5}
\det\!\left(1-t\mathfrak F_q\mid H^i_{\mathrm{DR}}(X/\mathbb F_q)\right)
=
\det\!\left(1-t\mathfrak F_q\mid H^i(X,\mathcal O_X)\right).
\end{equation}

%%%%%%%%%%  scan idx 419  |  folio 412  %%%%%%%%%%

\underline{Théorème} 3.4. \underline{Soit} $X$ \underline{un schéma séparé de
type fini sur} $\mathbb F_q$. \underline{Alors désignant par} $H_c^i$ \underline{la cohomologie
étale à supports propres} (SGA 4 XVII 5.1.9 (iii)), \underline{on a}

\begin{equation}\tag{3.4.1}
Z(t,X/\mathbb F_q)
\equiv
\prod_{i\geq0}
\det\!\left(1-t\varphi_q^{-1}\mid
H_c^i(\bar X,\mathbb Z/p\mathbb Z)\right)^{(-1)^{i+1}}
\mod p.
\end{equation}

Commençons par démontrer l'équivalence de (3.4.1) et (3.1.1).

Écrivons, pour le moment

\begin{equation}\tag{3.4.2}
Z(X)\quad\text{au lieu de }Z(t,X/\mathbb F_q)
\end{equation}

\begin{equation}\tag{3.4.3}
Z'(X)\quad\text{au lieu de }\quad
\prod\det\!\left(1-t\varphi_q^{-1}\mid
H_c^i(\bar X,\mathbb Z/p\mathbb Z)\right)^{(-1)^{i+1}}.
\end{equation}

\underline{Lemme} 3.5. \underline{Si} $\{U_i\}$ \underline{est un recouvrement fini de}
$X$ \underline{par des sous-schémas ouverts on a}

\begin{equation}\tag{3.5.1}
Z'(X)=
\prod_{j\geq0}
\prod_{1\leq i_0<\cdots<i_j\leq n}
Z'(U_{i_0}\cap\cdots\cap U_{i_j})^{(-1)^j}.
\end{equation}

\underline{Démonstration.} Résulte facilement de la suite spectrale de
localisation (SGA 4 XVII 6.2.10)

\[
E_1^{-a,b}=\bigoplus_{i_0<\cdots<i_b}
H_c^a(\bar U_{i_0}\cap\cdots\cap\bar U_{i_b},\mathbb Z/p\mathbb Z)
\Longrightarrow H_c^{-a+b}(\bar X,\mathbb Z/p\mathbb Z).
\]

On peut aussi l'obtenir par récurrence sur $n$, en appliquant 3.6 ci-dessous.

\underline{Lemme} 3.6. \underline{Si} $U$ \underline{est un ouvert de} $X$, $Y$ \underline{le fermé}
$X-U$,

\begin{equation}\tag{3.6.1}
Z'(X)=Z'(U)Z'(Y)
\end{equation}

\underline{Démonstration.} Résulte de la suite exacte de cohomologie

\begin{equation}\tag{3.6.2}
\longrightarrow H_c^i(\bar U,\mathbb Z/p\mathbb Z)
\longrightarrow H_c^i(\bar X,\mathbb Z/p\mathbb Z)
\longrightarrow H_c^i(\bar Y,\mathbb Z/p\mathbb Z)\longrightarrow.
\end{equation}

%%%%%%%%%%  scan idx 420  |  folio 413  %%%%%%%%%%

\underline{Lemme} 3.7 (Moebius). \underline{Soit} $B$ \underline{un ensemble fini},
$\{A_i\}$ \underline{un recouvrement fini par des sous-ensembles}. \underline{Alors}

\begin{equation}\tag{3.7.1}
\operatorname{card}(B)=
\sum_{j\geq0}(-1)^j
\sum_{i_0<\cdots<i_j}
\operatorname{card}(A_{i_0}\cap\cdots\cap A_{i_j}).
\end{equation}

\underline{Démonstration.} On a la suite spectrale de Leray pour la
cohomologie rationnelle de l'espace discret $B$ et le recouvrement ouvert
$\{A_i\}$

\[
E_1^{p,q}=\bigoplus_q
H^p(A_{i_0}\cap\cdots\cap A_{i_q},\mathbb Q)
\Longrightarrow H^{p+q}(B,\mathbb Q)
\]

et on calcule la caractéristique d'Euler-Poincaré de $B$.

\underline{Lemme} 3.8. \underline{Si} $\{U_i\}$ \underline{est un recouvrement fini de}
$X$ \underline{par des sous-schémas, alors}

\begin{equation}\tag{3.8.1}
Z(X)=
\prod_{j\geq0}
\prod_{i_0<\cdots<i_j}
Z(U_{i_0}\cap\cdots\cap U_{i_j})^{(-1)^j}.
\end{equation}

\underline{Démonstration.} On a la formule

\begin{equation}\tag{3.8.2}
\operatorname{card}(X(\mathbb F_{q^n}))
=\sum_{r\mid n}r.\operatorname{card}\{y\in X\mid\deg(y)=r\}
\end{equation}

qui entraîne

\begin{equation}\tag{3.8.3}
Z(X)=\exp\!\left(
\sum_{n\geq1}\operatorname{card}(X(\mathbb F_{q^n}))\frac{t^n}{n}
\right),
\end{equation}

et on conclut par 3.7.1.

\underline{Lemme} 3.9. \underline{Supposons que la formule} (3.4.1) (\underline{ou ce qui
revient au même}, (3.1.2) \underline{i}.\underline{e}. (3.1.1)) \underline{soit vraie pour} $X$ \underline{une hypersurface
projective}. \underline{Alors elle est vraie pour tout} $X$ \underline{séparé de type fini}.

%%%%%%%%%%  scan idx 421  |  folio 414  %%%%%%%%%%

\underline{Démonstration.} Par 3.8 et 3.6, (3.4.1) est vraie pour le
complémentaire d'une hypersurface, $\mathbb P^n-H$. Compte tenu de ce que $\bigcap(\mathbb P^n-H_i)=\mathbb P^n-\bigcup_iH_i$ est encore le complément d'une hypersurface, le recouvrement ouvert
$\{\mathbb P^n-H_i\}$ de $\mathbb P^n-\bigcap H_i$ permet d'appliquer 3.8 et
3.5 pour déduire (3.1.3) pour $\mathbb P^n-\bigcap H_i$. Par 3.8 et 3.6,
(3.4.1) est vrai pour $\bigcap H_i$, i.e. pour tout $X$ projectif. Une
nouvelle application de 3.8 et 3.5 donne 3.1.3 pour tout $X$ quasi-projectif,
en particulier pour tout $X$ affine. Finalement, pour $X$ séparé de type
fini, on prend un recouvrement fini par des ouverts affines, et on applique
3.8 et 3.5.

4. \underline{Le calcul de la fonction zêta d'une hypersurface à la Dwork}

4.0. Soit

\begin{equation}\tag{4.0.1}
f(X)=f(X_1,\ldots,X_{n+2})\in\mathbb F_q[X]
\quad\text{homogène de degré }d,
\end{equation}

dont l'annulation définit l'hypersurface $V\subset\mathbb P^{n+1}$, et
désignons par $V^*$ l'ouvert de $V$ où
$\prod_{i=1}^{n+2}X_i$ est inversible, $V^*_{\mathrm{aff}}$ le cône affine
de $V^*$.

Choisissons un caractère additif non-trivial

\begin{equation}\tag{4.0.2}
\chi:\mathbb F_{q^\nu}\longrightarrow\mu_p(C),
\end{equation}

$C$ étant un corps alg. clos fixé de car. nulle
(p.ex. $C=\mathbb C$), de sorte qu'on a

\begin{equation}\tag{4.0.3}
\text{pour }x\in\mathbb F_{q^\nu}\qquad
\sum_{y\in\mathbb F_{q^\nu}}\chi(xy)=
\begin{cases}
q^\nu&\text{si }\chi=0\\
0&\text{si }\chi\ne0.
\end{cases}
\end{equation}

Pour $\chi=(\chi_1,\ldots,\chi_{n+2})\in(\mathbb F_{q^\nu}^*)^{n+2}$, on a donc

%%%%%%%%%%  scan idx 422  |  folio 415  %%%%%%%%%%

\begin{equation}\tag{4.0.4}
1+\sum_{X_0\in\mathbb F_{q^\nu}^*}\chi(x_0f(x))
=\sum_{x_0\in\mathbb F_{q^\nu}}\chi(x_0f(x))
=
\begin{cases}
q^\nu&\text{si }f(x)=0\\
0&\text{sinon}.
\end{cases}
\end{equation}

En sommant (4.0.4) sur $x\in(\mathbb F_{q^\nu}^*)^{n+2}$, on obtient

\begin{equation}\tag{4.0.5}
\operatorname{card}V^*_{\mathrm{aff}}(\mathbb F_{q^\nu})
=\frac1{q^\nu}
\sum_{x\in(\mathbb F_{q^\nu}^*)^{n+2}}
\left\{1+\sum_{x_0\in\mathbb F_{q^\nu}^*}\chi(x_0f(x))\right\}
\end{equation}

et, compte tenu de ce que

\begin{equation}\tag{4.0.6}
\operatorname{card}V^*_{\mathrm{aff}}(\mathbb F_{q^\nu})
=(q^\nu-1)\operatorname{card}V^*(\mathbb F_{q^\nu})
\end{equation}

(4.0.5) se récrit

\begin{equation}\tag{4.0.7}
(q^\nu-1)\operatorname{card}V^*(\mathbb F_{q^\nu})
=\frac{(q^\nu-1)^{n+2}}{q^\nu}
+\frac1{q^\nu}
\sum_{(x_0,x)\in(\mathbb F_{q^\nu}^*)^{n+3}}\chi(x_0f(x)).
\end{equation}

Écrivons explicitement la forme $f(X)$ comme somme de monômes

\begin{equation}\tag{4.0.8}
f(X)=\sum a_wX^w,
\end{equation}

où $w=(w_1,\ldots,w_{n+2})$ avec $\sum w_i=d$,
$X^w\overset{\mathrm{dfn}}{=}X_1^{w_1}\cdots X_{n+2}^{w_{n+2}}$, de sorte
que (4.0.7) se récrit

\begin{equation}\tag{4.0.9}
(q^\nu-1)\operatorname{card}V^*(\mathbb F_{q^\nu})
=\frac{(q^\nu-1)^{n+2}}{q^\nu}
+\frac1{q^\nu}
\sum_{(x_0,x)\in(\mathbb F_{q^\nu}^*)^{n+3}}
\prod_w\chi(a_wx_0x^w).
\end{equation}

4.1. Il faut maintenant expliciter notre ``caractère additif à valeurs en
caractéristique zéro''. On a envie de prendre simplement

\begin{equation}\tag{4.1.0}
\chi:x\longmapsto
\exp\!\left(2\pi i\,\operatorname{tr}_{\mathbb F_{q^\nu}/\mathbb F_p}(x)\right.
=\xi_p^{\operatorname{tr}_{\mathbb F_{q^\nu}/\mathbb H_p}(x)}.
\end{equation}

%%%%%%%%%%  scan idx 423  |  folio 416  %%%%%%%%%%

Désignons par

\begin{align}
\tag{4.1.1}
\Omega&\quad\text{la clôture algébrique de }\mathbb Q_p\\
\tag{4.1.2}
\operatorname{ord}:\Omega^*&\longrightarrow\mathbb Q
\quad\text{l'ordinal, normalisé par }\operatorname{ord}(p)=1\\
\tag{4.1.3}
\xi_p&\in\Omega
\quad\text{une racine $p$-ième primitive de l'unité.}
\end{align}

On constate (cf. [2]) qu'il existe

\begin{equation}\tag{4.1.4}
\pi\in\mathbb Q_p(\xi_p),\qquad
\operatorname{ord}(\pi)=\frac1{p-1}
\end{equation}

tel que

\begin{equation}\tag{4.1.5}
\sum_{n\geq0}\frac{\pi^{p^n}}{p^n}=0,
\end{equation}

(NB. On constate aussitôt que le premier membre converge dans
$Q_p(\xi_p)$.)

On définit \underline{l'exponentielle d'Artin-Hasse}

\begin{equation}\tag{4.1.6}
E(X)=\exp\!\left(\sum_{n\geq0}\frac{X^{p^n}}{p^n}\right)
\in Q_p[[X]]
\end{equation}

qui, oh miracle, se trouve dans $\mathbb Z_p[[X]]$. On définit

\begin{equation}\tag{4.1.7}
\Theta(X)=E(\pi X),
\end{equation}

dont les propriétés fondamentales sont résumées dans le lemme suivant :

\underline{Lemme} 4.2.

\begin{equation}\tag{4.2.1}
\Theta(X)=\sum B_nX^n,\qquad B_0=1\quad A_1=\pi,\quad
\operatorname{ord}(B_n)\geq\frac n{p-1}.
\end{equation}

\begin{equation}\tag{4.2.2}
\Theta(1)=\xi_p,\quad
\text{\uline{une racine $p$'ième primitive de l'unité.}}
\end{equation}

(4.2.3) \underline{Soit $t\in\Omega$, et supposons que, pour un entier
$a>0$ donné, $t=t^{p^a}$ (i.e. $t$ est un ``représentant de Teichmüller''
dans $\mathbb W(\mathbb F_{q^a})\subset\Omega$). Alors
$\sum_{j=0}^{a-1}t^{p^j}\in\mathbb Z_p$, de sorte qu'on peut définir la
série formelle}

%%%%%%%%%%  scan idx 424  |  folio 417  %%%%%%%%%%

\[
(\Theta(X))^{\sum_{j=0}^{a-1}t^{p^j}},
\]

et on a une identité de séries formelles

\begin{equation}\tag{4.2.3.1}
\prod_{j=0}^{a-1}\Theta(t^{p^j}X)
=(\Theta(X))^{\sum_{j=0}^{a-1}t^j}.
\end{equation}

\underline{Démonstration.} Compte tenu de ce qui
$E(X)\in\mathbb Z_p[[X]]$, (4.2.1) se déduit de la définition. Quant à
(4.2.2), on a

\begin{equation}\tag{4.2.4}
\log(\Theta(1))=0
\qquad\text{par le choix de $\pi$ (4.1.5)}
\end{equation}

et

\begin{equation}\tag{4.2.5}
\operatorname{ord}(\Theta(1)-1)=\frac1{p-1}
\qquad\text{par (4.2.1).}
\end{equation}

Par (4.2.4), il existe un entier $n\geq0$ tel que
$\Theta(1)^{p^n}=1$, (c'est la caractérisation du noyau du logarithme) et
(4.2.5) entraîne que $n=1$.

Quant à (4.2.3.1), en prenant les logarithmes, il résulte d'un calcul évident,
une fois prouvé que l'on a bien
$\sum_{j=0}^{a-1}t^{p^j}\in\mathbb Z_p$. Or on voit aussitôt que les
$t^{p^j}$ ($0\leq j\leq a-1$) sont les conjugués de
$t\in\mathbb W(\mathbb F_{q^a})$ sur $\mathbb Z_p$, donc

\[
\sum_{j=0}^{a-1}t^{p^j}
=\operatorname{Tr}_{\mathbb W(\mathbb F_{q^a})/\mathbb Z_p}(t)\in\mathbb Z_p.
\]

En prenant la valeur en $X=1$ des deux membres de (4.2.3.1) nous trouvons le

\underline{Corollaire} 4.3. \underline{Soit $t\in\Omega$ vérifiant $t^{p^a}=t$.
Désignons par $\bar t\in\mathbb F_{p^a}$ sa classe résiduelle. Alors on a}

\begin{equation}\tag{4.3.1}
\prod_{j=0}^{a-1}\Theta(t^{p^j})
=\xi_p^{\operatorname{tr}_{\mathbb F_{p^a}/\mathbb F_p}(\bar t)}.
\end{equation}

%%%%%%%%%%  scan idx 425  |  folio 418  %%%%%%%%%%

4.4. Rappelons que le corps résiduel de l'anneau des entiers de
$\mathbb Q_p(\xi_p)$ est $\mathbb F_p$. Soient

\begin{equation}\tag{4.4.0}
K=\text{l'extension non-ramifiée de $\mathbb Q_p(\xi_p)$ à corps résiduel
$\mathbb F_q$},
\end{equation}

\begin{equation}\tag{4.4.1}
\tau\in\Gal(K/Q_p(\xi_p))
\end{equation}

le "Frobenius", i.e. l'automorphisme se réduisant suivant l'automorphisme
de Frobenius de $\mathbb F_q/\mathbb F_p$, $x\longmapsto x^p$,

\begin{equation}\tag{4.4.2}
\text{Soit }F(X)=\sum A_WX^W
\text{ le relèvement en $K[X_1,\ldots,X_{n+2}]$ de $f(X)$}
\end{equation}

dont les coefficients vérifient $A_W^q=A_W$. Mettant ensemble (4.3.1) et
(4.0.9), on trouve (compte tenu de la définition (4.1.0) de $X$) :

\begin{equation}\tag{4.4.3}
\left\{
\begin{aligned}
&(q^\nu-1)\operatorname{card}V^*(\mathbb F_{q^\nu})
-\frac{(q^\nu-1)^{n+2}}{q^\nu}\\
&\quad=\frac1{q^\nu}
\sum_{(X_0,X)\in(\mu_{q^\nu-1}(\Omega))^{n+3}}
\prod_W\prod_{j=0}^{\operatorname{ord}(q^\nu)-1}
\Theta\!\left(A_W^{p^j}X_0^{p^j}X^{p^jW}\right).
\end{aligned}
\right.
\end{equation}

Définissons

\[
H(X)=H(X_0,\ldots,X_{n+2})\in K[[X_0,\ldots,X_{n+2}]]
\]

par

\begin{equation}\tag{4.4.4}
H(X)=\prod_W
\Theta\!\left(A_WX_0X_1^{W_1}\cdots X_{n+2}^{W_{n+2}}\right).
\end{equation}

Nous faisons agir $\Gal(K/\mathbb Q_p(\xi_p))$ sur
$K[[X_0,\ldots,X_{n+2}]]$ à travers son action sur les coefficients ;
l'action composante de $\sigma\in\Gal(K/\mathbb Q_p(\xi_p))$ sera notée

\begin{equation}\tag{4.4.5}
G(X)\longmapsto G^\sigma(X).
\end{equation}

%%%%%%%%%%  scan idx 426  |  folio 419  %%%%%%%%%%

On a en particulier

\[
H^{\tau^j}(X)=\prod_W
\Theta\!\left(A_W^{\tau^j}X_0X_1^{W}\cdots
X_{n+2}^{W_{n+2}}\right).
\]

Compte tenu de ce que $A_W^{\tau^j}=A_W^{p^j}$, (4.4.3) se récrit alors

\begin{equation}\tag{4.4.6}
\begin{aligned}
(q^\nu-1)\operatorname{Card}V^*(\mathbb F_{q^\nu})
&=\frac{(q^\nu-1)^{n+2}}{q^\nu}\\
&\quad+\frac1{q^\nu}
\sum_{(X_0,X)\in(\mu_{q^\nu-1}(\Omega))^{n+3}}
\prod_{j=0}^{\operatorname{ord}(q^\nu)-1}H^{\tau^j}(X^{p^j})
\end{aligned}
\end{equation}

4.5. Nous allons maintenant exprimer la somme dans (4.4.6) comme une trace.

Pour chaque $b\in\mathbb Q$, on définit une $K$-algèbre de Banach $p$-adique

\begin{equation}\tag{4.5.1}
L(b)=\text{le sous-anneau de $K[[X_0,\ldots,X_{n+2}]]$ formé des séries}
\end{equation}

qui s'écrivent

\begin{equation}\tag{4.5.2}
\sum_{dU_0=\sum_{i\geq1}U_i}
B_UX_0^{U_0}X_1^{U_1}\cdots X_{n+2}^{U_{n+2}}
\qquad (U=(U_0,\ldots,U_{n+2}))
\end{equation}

et qui vérifient

\begin{equation}\tag{4.5.3}
\inf_U\bigl(\operatorname{ord}(B_U)-bU_0\bigr)>-\infty.
\end{equation}

C'est une algèbre de Banach $p$-adique pour la fonction d'ordinal
(l'opposé du logarithme de la norme)

\begin{equation}\tag{4.5.4}
\inf_U\bigl(\operatorname{ord}(B_U)-bU_0\bigr)
=\operatorname{ord}\!\left(\sum B_UX^U\right).
\end{equation}

On vérifie sur 4.4.4 et (4.2.1) que l'on a

\[
H(X)\in L\!\left(\frac1{p-1}\right).
\]

%%%%%%%%%%  scan idx 427  |  folio 420  %%%%%%%%%%

Introduisons des opérateurs.

Pour chaque entier $k>0$, on introduit un opérateur continu d'espaces de
Banach sur $K$

\begin{equation}\tag{4.5.5}
\left\{
\begin{aligned}
\psi_k:L(b)&\longrightarrow L(kb)\\
\psi_k\!\left(\sum B_UX^U\right)&=\sum B_{kU}X^U.
\end{aligned}
\right.
\end{equation}

Pour chaque couple $(b,b')$, $b'<b$, il y a un
\underline{opérateur de} "\underline{restriction}" qui est complètement continu
(i.e. une limite uniforme des opérateurs de rang fini)

\begin{equation}\tag{4.5.6}
\left\{
\begin{aligned}
L(b)&\xrightarrow{\ \text{"restriction"}\ }L(b')\\
\sum B_UX^U&\longmapsto\sum B_UX^U.
\end{aligned}
\right.
\end{equation}

Chaque élément $G$ de l'algèbre $L(b)$ donne l'endomorphisme continu de la
multiplication par $G$

\begin{equation}\tag{4.5.7}
\begin{aligned}
G:L(b)&\longrightarrow L(b)\\
\eta&\longmapsto\eta G.
\end{aligned}
\end{equation}

Ces opérateurs sont liés par le lemme suivant, dont la démonstration se
trouve dans [2] (cf. aussi [4]).

%%%%%%%%%%  scan idx 428  |  folio 421  %%%%%%%%%%

\underline{Lemme} 4.5.8 (\underline{ou lemme de la trace} (DWORK)).
\underline{Soient} $b\in\mathbb Q$, $b>0$, $k\in\mathbb Z$, $k\geq2$, \underline{et}
$G\in L(b/k)$. \underline{Alors le composé}

\[
L(b)\xrightarrow{\ \text{"restriction"}\ }L(b/k)
\xrightarrow{\ \times G\ }L(b/k)
\xrightarrow{\ \psi_k\ }L(b),
\]

\underline{noté} $\psi_k\circ G$, \underline{est complètement continu, et sa trace}
[cf. [4]] \underline{est donnée par la formule}

\begin{equation}\tag{4.5.9}
\operatorname{tr}(\psi_k\circ G)
=\frac1{(k-1)^{n+2}}
\sum_{X\in(\mu_{k-1}(\Omega))^{n+3}}G(X).
\end{equation}

Appliquant ce lemme au cas

\[
\left\{
\begin{aligned}
k&=q^\nu\\
G(X)&=\prod_{j=0}^{\operatorname{ord}(q^\nu)-1}
H^{\tau^j}(X^{p^j})
\in L\!\left(\frac{p}{(p-1)q^\nu}\right)\\
b&=\frac1{p-1}
\end{aligned}
\right.
\]

(4.4.6) se récrit

\begin{equation}\tag{4.5.10}
\begin{aligned}
(q-1)\operatorname{card}V^*(\mathbb F_{q^\nu})
&=\frac{(q^\nu-1)^{n+2}}{q^\nu}\\
&\quad+\frac{(q^\nu-1)^{n+3}}{q^\nu}
\operatorname{tr}\!\left(
\left.\psi_q
\prod_{j=0}^{\operatorname{ord}(q^\nu)-1}H^{\tau^j}(X^{p^j})
\right|L\!\left(\frac1{p-1}\right)\right).
\end{aligned}
\end{equation}

Soit $L^0(b)$ l'idéal des éléments de $L(b)$ qui "s'annulent en $0$",
i.e. "sans terme constant". Il est stable sous $\psi\circ G$, et on a
évidemment

\begin{equation}\tag{4.5.11}
\operatorname{tr}(\psi_k\circ G)
=G(0)+\operatorname{tr}\!\left(
\left.\psi_k\circ G\right|L^0\right).
\end{equation}

Compte tenu de ce que $H(0)=1$, (4.5.10) se récrit

\begin{equation}\tag{4.5.12}
\begin{aligned}
\operatorname{card}V^*(\mathbb F_{q^\nu})
&=(q^\nu-1)^{n+1}\\
&\quad+\frac{(q^\nu-1)^{n+2}}{q}
\operatorname{tr}\!\left(
\left.\psi_q\circ
\prod_{j=0}^{\operatorname{ord}(q^\nu)-1}H^{\tau^j}(X^{p^j})
\right|L^0\!\left(\frac1{p-1}\right)\right).
\end{aligned}
\end{equation}

%%%%%%%%%%  scan idx 429  |  folio 422  %%%%%%%%%%

Considérons maintenant l'opérateur $\tau^{-1}$-linéaire

\begin{equation}\tag{4.5.13}
\tau^{-1}\circ\psi_p\circ H(X):L(b)\longrightarrow L(b).
\end{equation}

On constate que ses itérés sont donnés par

\begin{equation}\tag{5.4.14}
(\tau^{-1}\circ\psi_p\circ H(X))^a
=\tau^{-a}\circ\psi_{p^a}\circ
H(X)H^\tau(X^p)\cdots H^{\tau^{a-1}}(X^{p^{a-1}}),
\end{equation}

et on note que, si $\operatorname{ord}(q)\mid a$, on a
$\tau^{-a}=\operatorname{id}$. Ceci posé, on définit une opération
$\tau^{-1}$-linéaire

\begin{equation}\tag{4.5.15}
\left\{
\begin{aligned}
\alpha:L^0\!\left(\frac1{p-1}\right)
&\longrightarrow L^0\!\left(\frac1{p-1}\right)\\
\alpha&=\frac1p\,\tau^{-1}\circ\psi_p\circ H(X),
\end{aligned}
\right.
\end{equation}

de sorte que (4.5.12) se récrit

\begin{equation}\tag{4.5.16}
\operatorname{card}V^*(\mathbb F_{q^\nu})
=(q^\nu-1)^{n+1}
+(q^\nu-1)^{n+2}
\operatorname{tr}\!\left(
\left.(\alpha)^{\operatorname{ord}(q^\nu)}
\right|L^0\!\left(\frac1{p-1}\right)\right..
\end{equation}

Désignons par

\begin{equation}\tag{4.5.17}
S=\text{l'ensemble }\{1,\ldots,n+2\}.
\end{equation}

Pour tout sous-ensemble $A\subset S$, $A\ne\varnothing$, on définit un
sous-espace de $L(b)$ :

\begin{equation}\tag{4.5.18}
\begin{aligned}
L^A(b)={}&\text{ensemble des séries $\sum B_UX^U\in L(b)$ telles que, pour chaque
monôme $X^U$ avec $B_U\ne0$,}\\
&\text{les $U_i>0$ sont exactement $U_0$ et les $U_j$, $j\in A$}.
\end{aligned}
\end{equation}

On a

\begin{equation}\tag{4.5.19}
L^0(b)=\bigoplus_{\substack{A\subset S\\A\ne\varnothing}}L^A(b)
\end{equation}

%%%%%%%%%%  scan idx 430  |  folio 423  %%%%%%%%%%

et, pour chaque sous-ensemble

\[
B\subset S,\qquad B\ne\varnothing,
\]

on a

\begin{equation}\tag{4.5.20}
\bigoplus_{A\supset B}L^A(b)
=L(b)\cap\text{ l'idéal }
\left(\prod_{i\in B}X_i\right)K[[X]].
\end{equation}

Évidemment chaque sous-espace $\bigoplus_{A\subseteq B}L^A(b)$ est stable
par $\alpha$ ; il en résulte qu'il en est de même pour
$\bigoplus_{A\supsetneq B}L^A(b)$, qui est somme de sous-espaces du type
précédent.

On définit

\begin{equation}\tag{4.5.21}
\alpha_B=\text{l'endomorphisme $\tau^{-1}$-linéaire induit par $\alpha$ sur}
\quad
\frac{\displaystyle\bigoplus_{A\supset B}
L^A\!\left(\frac1{p-1}\right)}
{\displaystyle\bigoplus_{A\supsetneq B}
L^A\!\left(\frac1{p-1}\right)}
\simeq L^B\!\left(\frac1{p-1}\right).
\end{equation}

Explicitement, on obtient $\alpha_B$ par la formule

\begin{equation}\tag{4.5.22}
\alpha_B=\frac{\tau^{-1}}p\,\psi_p\,H_B(X),
\end{equation}

opération sous laquelle $L^B$ est stable, où $H_B(X)$ se déduit de $H(X)$
en y faisant la substitution $X_i=0$ pour $i\notin\{0\}\cup B$.

Récrivons encore (4.5.16)

\begin{equation}\tag{4.5.2.3}
\operatorname{card}V^*(\mathbb F_{q^\nu})
=(q^\nu-1)^{n+1}
+(q^\nu-1)^{n+2}
\sum_{\substack{A\subset S\\A\ne\varnothing}}
\operatorname{tr}\!\left((\alpha_A)^{\operatorname{ord}(q^\nu)}\right).
\end{equation}

Définissons, pour $A\subset S$, $A\ne\varnothing$, l'hypersurface

\begin{equation}\tag{4.5.24}
V_A\subset\mathbb P^{\operatorname{card}(A)-1}
\simeq
\mathbb P^{n+1}\cap
\{X_i=0\ \forall i\notin A\}
\end{equation}

définie par $f_A(X)$, le polynôme déduit de $f(X)$ y faisant la
substitution en $X_i=0$ pour $i\notin A$.

%%%%%%%%%%  scan idx 431  |  folio 424  %%%%%%%%%%

Compte tenu de la compatibilité de la formation de la série $H$ à partir de
$f(X)$ à l'opération de substitution de 0 pour certaines des variables
(4.5.23) donne, pour chaque $A\subseteq S$, $A\ne\varnothing$

\begin{equation}\tag{4.5.25}
\operatorname{card}(V_A^*(\mathbb F_{q^\nu}))
=(q^\nu-1)^{\operatorname{car}(A)-1}
+(q^\nu-1)^{\operatorname{card}(A)}
\sum_{B\subset A}
\operatorname{tr}\!\left((\alpha_B))^{\operatorname{ord}(q^\nu)}\right).
\end{equation}

D'autre part, on a

\begin{equation}\tag{4.5.26}
\operatorname{card}V(\mathbb F_{q^\nu})
=\sum_{\substack{A\subset S\\A\ne\varnothing}}
\operatorname{card}V_A^*(\mathbb F_{q^\nu}).
\end{equation}

Après une sommation facile, (4.5.25) donne alors

\begin{equation}\tag{4.5.27}
\begin{aligned}
\operatorname{card}V(\mathbb F_{q^\nu})
={}&1+q^\nu+q^{2\nu}+\cdots
+q^{\nu(\operatorname{card}(S)-1)}\\
&+\sum_{\substack{B\subseteq S\\B\ne\varnothing}}
(q^\nu-1)^{\operatorname{card}(B)}
q^{\nu\operatorname{card}(S-B)}
\operatorname{tr}\!\left((\alpha_B)^{\operatorname{ord}(q^\nu)}\right).
\end{aligned}
\end{equation}

Notons que $\alpha$ donc les $\alpha_B$ étant $\tau^{-1}$ linéaire, les
$\alpha_B^{\operatorname{ord}(q)}$ sont $K$-linéaires. Définissons alors les
séries caractéristiques

\[
\Delta_B(t)\in1+tK[[t]]
\]

\begin{equation}\tag{4.5.28}
\Delta_B(t)=\det\!\left(
1-t\alpha_B^{\operatorname{ord}(q)}
\mathrel{\Big|}L^B\!\left(\frac1{p-1}\right)\right).
\end{equation}

Introduisons des endomorphismes de $1+tK[[t]]$

\begin{equation}\tag{4.5.29}
\varphi:g(t)\longmapsto g^\varphi(t)=g(qt)
\end{equation}

et

\begin{equation}\tag{4.5.30}
\delta=1-\varphi:g(t)\longmapsto
g^\delta(t)=g(t)/g(qt)
\end{equation}

%%%%%%%%%%  scan idx 432  |  folio 425  %%%%%%%%%%

Compte tenu de ce que

\begin{equation}\tag{4.5.31}
Z(t,V/\mathbb F_q)
=\exp\!\left(
\sum_{\nu\geq1}\operatorname{card}(V(\mathbb F_{q^\nu}))
\frac{t^\nu}{\nu}\right)
\end{equation}

et

\begin{equation}\tag{4.5.32}
\det\!\left(1-t\alpha_B^{\operatorname{ord}(q)}\right)
=\exp\!\left(
-\sum_{\nu\geq1}
\operatorname{tr}\!\left((\alpha_B)^{\nu\operatorname{ord}q}\right)
\frac{t^\nu}{\nu}\right)
\end{equation}

(4.5.27) se récrit

\begin{equation}\tag{4.5.23}
Z(t,V/\mathbb F_q)\prod_{i=0}^{n+1}(1-q^it)
=\prod_{j=1}^{n+2}
\left(
\prod_{\substack{B\subset S\\\operatorname{card}(B)=j}}
\Delta_B(t)
\right)^{-(-\delta)^j\varphi^{n+2-j}}.
\end{equation}

4.6. Nous allons étudier maintenant l'intégralité des séries $\Delta_B$.
Posons

\begin{equation}\tag{4.6.1}
\mathcal O=\text{l'anneau des entiers de $K$}
\end{equation}

et définissons une \underline{filtration} sur $1+t\mathcal O[[t]]$ par

\begin{equation}\tag{4.6.2}
F^i(1+t\mathcal O[[t]])
=1+q^it\mathcal O[[q^it]]
=\varphi^i(1+t\mathcal O[[t]]).
\end{equation}

On a

\begin{align}
\varphi(F^i)&\subset F^{i+1} \tag{4.6.3}\\
\delta(F^i)&\subset F^i \tag{4.6.4}\\
\delta&\text{ induit l'identité sur }F^i/F^{i+1}. \tag{4.6.5}
\end{align}

\underline{Lemme} 4.6.6.
$\Delta_B(t)\in1+t\mathcal O[[t]]$, \underline{et si un entier} $j$ \underline{vérifie}

\begin{equation}\tag{4.6.7}
jd<\operatorname{card}(B)
\end{equation}

\underline{alors on a}

\begin{equation}\tag{4.6.8}
\Delta_B(t)\in F^j(1+t\mathcal O[[t]]).
\end{equation}

%%%%%%%%%%  scan idx 433  |  folio 426  %%%%%%%%%%

\underline{Démonstration.}
Les espaces de Banach
$L^B\!\left(\frac1{p-1}\right)$ sont tous \og du type $B(I)$\fg{} au sens de
Serre [7], i.e. ils admettent des \og bases orthonormales\fg{}, à savoir, les
monômes

\begin{equation}\tag{4.6.9}
\pi^{U_0}X_0^{U_0}X_1^{U_1}\cdots X_{n+2}^{U_{n+2}}
\end{equation}

\noindent
qui y appartiennent. Que c' une base \og orthonormale\fg{} veut dire par
définition qu'on a

\begin{equation}\tag{4.6.10}
\operatorname{ord}\!\left(
\sum C_U\pi^{U_0}X_0^{U_0}\cdots X_{n+2}^{U_{n+2}}
\right)
=\min_U\operatorname{ord}(C_U).
\end{equation}

Soit

\begin{equation}\tag{4.6.11}
\mathbf R_B=\text{le $\mathcal O$-sous-module de
$L^B\!\left(\frac1{p-1}\right)$ formé des éléments d'ordinal $\geq0$}.
\end{equation}
Pour démontrer le lemme, il suffit de démontrer

\begin{equation}\tag{4.6.12}
\text{si }jd<\operatorname{card}(B),\text{ alors }
(\alpha_B)^{\operatorname{ord}(q)}R_B\subset q^jR_B
\end{equation}

ou, a fortiori, de démontrer

\begin{equation}\tag{4.6.13}
\text{si }jd<\operatorname{card}(B),\text{ alors }
\alpha_B(R_B)\subset p^jR_B.
\end{equation}

Or, on a (4.5.22)

\begin{equation}\tag{4.6.14}
\alpha_B=\frac{\tau^{-1}}p\,\psi_p\circ H_B(X).
\end{equation}

Compte tenu de ce que

\begin{equation}\tag{4.6.15}
H(X)\in L\!\left(\frac1{p-1}\right),
\qquad \operatorname{ord}(H(X))\geq0
\end{equation}

on a

\begin{equation}\tag{4.6.16}
H_B(X)\cdot R_B\subset R_B
\end{equation}

et évidemment

\begin{equation}\tag{4.6.17}
\tau^{-1}(R_B)=R_B,
\end{equation}

%%%%%%%%%%  scan idx 434  |  folio 427  %%%%%%%%%%

de sorte que (4.6.13) sera conséquence de

\begin{equation}\tag{4.6.18}
\text{si }jd<\operatorname{card}(B),\text{ alors }
\psi_p(R_B)\subset p^{j+1}R_B.
\end{equation}

Vérifions (4.6.18) monôme par monôme. On a

\begin{equation}\tag{4.6.19}
\begin{cases}
\psi_p\!\left(\pi^{U_0}X^{U_0}\cdots X_{n+2}^{U_{n+2}}\right)=0,
&\text{sauf si }\forall i\quad p\mid U_i\\
\text{et}\\[3pt]
\psi_p\!\left(
\pi^{pU_0}X_0^{pU_0}\cdots X_{n+2}^{pU_{n+2}}
\right)
=\pi^{(p-1)U_0}\cdot
\underbrace{\pi^{U_0}X_0^{U_0}\cdots X_{n+2}^{U_{n+2}}}_{\in R_B}.
\end{cases}
\end{equation}

Or, les conditions pour que l'on ait
$X_0^{U_0}\cdots X_{n+2}^{U_{n+2}}
\in L^B\!\left(\frac1{p-1}\right)$ sont

\begin{equation}\tag{4.6.20}
\begin{cases}
dU_0=\displaystyle\sum_{i\geq1}U_i,\\
\text{pour }i\geq1,\text{ on a }U_i>0\Longleftrightarrow i\in B
\end{cases}
\end{equation}

d'où, en particulier,

\begin{equation}\tag{4.6.21}
dU_0\geq\operatorname{card}(B).
\end{equation}

On a donc

\begin{equation}\tag{4.6.22}
\operatorname{ord}\!\left(\pi^{(p-1)U_0}\right)
=U_0\geq\frac{\operatorname{card}(B)}d>j
\end{equation}

et, comme $U_0\in\mathbb Z$,

\begin{equation}\tag{4.6.23}
\operatorname{ord}\!\left(\pi^{(p-1)U_0}\right)
=U_0\geq j+1.
\end{equation}

Mettant ceci ensemble avec (4.6.19), on trouve (4.6.18).

%%%%%%%%%%  scan idx 435  |  folio 428  %%%%%%%%%%

\section*{5. \underline{Un théorème d'AX} [1], (cf. aussi [4], section 7).}

\underline{Théorème 5.0.}
\underline{Soient} $V$ \underline{une hypersurface dans} $\mathbb P^{n+1}$ \underline{de degré} $d>0$, \underline{définie
sur} $\mathbb F_q$, \underline{et} $j$ \underline{un entier} : \underline{si on a}

\begin{equation}\tag{5.0.1}
jd<n+2,
\end{equation}

\underline{alors}

\begin{equation}\tag{5.0.2}
Z(t,V/\mathbb F_q)\prod_{i=0}^{n+1}(1-q^it)
\in1+q^jt\mathbb Z[[q^jt]].
\end{equation}

et l'inverse de tout zéro ou pôle de

\begin{equation}\tag{5.0.3}
Z(t,V/\mathbb F_q)\prod_{i=0}^{n+1}(1-q^it)
\end{equation}

\noindent
\underline{est de la forme} $q^j\cdot(\text{\underline{un entier algébrique}})$.

\underline{Théorème 5.0.bis.}
\underline{Soit} $V_{\mathrm{aff}}$ \underline{une hypersurface dans} $\mathbb A^{n+2}$ \underline{définie
par une forme homogène de degré} $d$, \underline{définie sur} $\mathbb F_q$. \underline{Si on a}

\begin{equation}\tag{5.0.7}
jd<n+2
\end{equation}

\underline{alors}

\begin{equation}\tag{5.0.5}
Z(t,V_{\mathrm{aff}}/\mathbb F_q)
\in1+q^jt\mathbb Z[[q^jt]].
\end{equation}

\underline{En particulier},

\begin{equation}\tag{5.0.6}
\operatorname{card}\bigl(V_{\mathrm{aff}}(\mathbb F_q)\bigr)
\equiv0\mod{q^j},
\end{equation}

\underline{et par suite}

\begin{equation}\tag{5.0.6 bis}
\operatorname{card}\bigl(V(\mathbb F_q)\bigr)
\geq\frac{q^j-1}{q-1}>0
\end{equation}

\noindent
(\underline{ce qui améliore le th. de CHEVALLEY--WARNING} [8]).

%%%%%%%%%%  scan idx 436  |  folio 429  %%%%%%%%%%

\underline{Démonstration.}
L'implication (5.0.5) $\Rightarrow$ (5.0.6) est claire, en prenant le
coefficient de $t$ dans la fonction zêta. Les assertions (5.0.4) et (5.0.2)
sont équivalentes, compte tenu de ce que

\begin{equation}\tag{5.0.7}
\operatorname{card}V_{\mathrm{aff}}(\mathbb F_{q^\nu})
=1+(q^\nu-1)\operatorname{card}V(\mathbb F_{q^\nu}),
\end{equation}

d'où

\begin{equation}\tag{5.0.8}
Z(t,V_{\mathrm{aff}}/\mathbb F_q)
=\frac{Z(t,V/\mathbb F_q)}
{Z(qt,V/\mathbb F_q)(1-t)}
\end{equation}

et

\begin{equation}\tag{5.0.9}
Z(t,V/\mathbb F_q)
=\prod_{i\geq0}\left[(1-q^if)
Z(q^it,V_{\mathrm{aff}}/\mathbb F_q)\right].
\end{equation}

Les énoncés (5.0.2) et (5.0.3) sont équivalents encore, par le théorème de
Fatou [3], grâce à la rationalité de la fonction zêta. Comme on a

\begin{equation}\tag{5.0.10}
Z(t,V/\mathbb F_q)\in\mathbb Z[[t]],
\end{equation}

il suffira de démontrer

\begin{equation}\tag{5.0.11}
Z(t,V/\mathbb F_q)\prod_{i=0}^{n+1}(1-q^it)
\in F^j(1+t\mathcal O[[t]]).
\end{equation}

Compte tenu de (4.5.33) et (4.6.4)--(4.6.5), il suffit de démontrer, pour
$B\subset\{1,\ldots,n+2\}$, $B\ne\varnothing$, qu'on a

\begin{equation}\tag{5.0.12}
\Lambda_B(t)\in F^{j+\operatorname{card}(B)-(n+2)}.
\end{equation}

Nous pourrons terminer la démonstration en appliquant le lemme (4.4.6), une
fois démontrée l'assertion :

\begin{equation}\tag{5.0.13}
\text{si }jd<n+2,\text{ alors }
\bigl(j+\operatorname{card}(B)-(n+2)\bigr)d
<\operatorname{card}(B),
\end{equation}

qui se récrit

%%%%%%%%%%  scan idx 437  |  folio 430  %%%%%%%%%%

\begin{equation}\tag{5.0.14}
\text{si }jd<n+2,\text{ alors }
(j-(n+2))d<\operatorname{card}(B)(1-d).
\end{equation}

Mais

\[
\inf_{\substack{B\subset\{1,\ldots,n+2\}\\B\ne\varnothing}}
\operatorname{card}(B)(1-d)
=(n+2)(1-d),
\]

de sorte qu'il suffit de démontrer

\begin{equation}\tag{5.0.15}
\text{si }jd<n+2,\text{ alors }
(j-(n+2))d<(n+2)(1-d),
\end{equation}

ce qui est vrai.

\underline{Remarque 5.1.}
La condition

\begin{equation}\tag{5.1.1}
jd<n+2
\end{equation}

admet une interprétation cohomologique. Soit $V$ une hypersurface \underline{lisse} de
degré $d$ et de dimension $n$. Posons, comme dans XI 2.1, pour $a+b=n$

\begin{equation}\tag{5.1.2}
h_0^{a,b}
=\dim H^b(V,\underline{\Omega}_V^a)-\delta_{a,b}
=\dim E^{a,b}(V)
\end{equation}

où $E^{a,b}(V)$ est la composante de bidegré $(a,b)$ de la \og cohomologie
évanescente de Hodge\fg{} (comparer XVII 5.4.4)

\begin{equation}\tag{5.1.3}
E^{a,b}(V)
=\operatorname{Coker}\!\left(
H^b(\mathbb P^{n+1},\underline{\Omega}_{\mathbb P^{n+1}}^a)
\hookrightarrow H^b(V,\underline{\Omega}_V^a)\right..
\end{equation}

Alors, par XI, 2.8,

\begin{equation}\tag{5.1.4}
jd<n+2
\Longleftrightarrow
\left[h_0^{a,n-a}=0\text{ pour }a<j\text{ et }a>n-j\right],
\end{equation}

i.e. la condition signifie que le \og niveau de Hodge\fg{} (XI 2.7) de la
\og cohomologie de Hodge évanescente\fg{} de $V$ :

%%%%%%%%%%  scan idx 438  |  folio 431  %%%%%%%%%%
\begin{equation*}
E_{\mathrm{DR}}^n(V)
=\operatorname{Coker}\!\left(
H_{\mathrm{Hdg}}^n(\mathbb P^{n+1})
\hookrightarrow H_{\mathrm{Hdg}}^n(V)\right),
\end{equation*}


est $\leq n-2j$, i.e. son co-niveau est $\geq j$. D'autre part, la
conclusion signifie aussi que les valeurs propres du Frobenius
$\mathfrak F_q$ opérant sur la cohomologie évanescente $\ell$-adique
($\ell\ne p$) sont des entiers algébriques (XXI A.2) \underline{divisibles par} $q^j$.
Le résultat de AX 5.0. vient donc à l'appui de la conjecture suivante,
suggérée par le yoga du niveau des motifs, et qui met en relations les
classiques conjectures de HODGE [5] et de TATE [5] :

\underline{Conjecture} 5.2.
\underline{Soient} $f:X\to S$ \underline{un morphisme propre et lisse}, \underline{avec} $S$ \underline{intègre}, \underline{de type
fini sur} $\operatorname{Spec}(\mathbb Z)$ \underline{et dominant}
$\operatorname{Spec}(\mathbb Z)$, $i,j$ \underline{des entiers}. \underline{Supposons que le
co-niveau de Hodge du} $H^i$ \underline{de la fibre générique de} $f$ \underline{soit} $\geq j$.
\underline{Alors pour tout point} $s\in S$, \underline{le coniveau} (XX 2.0.6) \underline{de}
$H^i(X_s,\mathbb Q_\ell)$ \underline{est} $\geq j$, \underline{donc} (XX 2.2) \underline{si} $s$ \underline{est un point
fermé de} $S$, \underline{les valeurs propres de Frobenius opérant sur}
$H^i(X_s,\mathbb Q_\ell)$ \underline{sont des entiers algébriques divisibles par}
$q^j$, \underline{où} $q=\operatorname{card}(k(s))$.

\section*{6. Fin de la démonstration de la formule de congruence.}

Compte tenu de (4.6.6) et (4.6.3)--(4.6.4), (4.5.24) donne une formule de
congruence

\begin{equation}\tag{6.0.1}
Z(t,V/\mathbb F_q)(1-t)
\equiv\Delta_S(t)^{(-1)^{n+1}}
\mod{1+qt\mathcal O[[t]]}.
\end{equation}

Passons au calcul de $\Delta_S(t)$ modulo $\pi\mathcal O[[t]]$. Définissons
(cf. (4.6.11) :

\begin{equation}\tag{6.0.2}
W_S=R_B\otimes_{\mathcal O}\mathbb F_q,
\end{equation}

%%%%%%%%%%  scan idx 439  |  folio 432  %%%%%%%%%%

\begin{equation}\tag{6.0.3}
\bar\alpha_S:W_S\longrightarrow W_S
\end{equation}

\noindent
l'opération $\tau^{-1}$-linéaire induite par $\alpha_S$,

\begin{equation}\tag{6.0.4}
R_S(t)\text{ le sous-module de $R_S$ forme des éléments de degré $i$ en $X_0$},
\end{equation}

\begin{equation}\tag{6.0.5}
W_S(i)=R_S(i)\otimes_{\mathcal O}\mathbb F_q.
\end{equation}

\underline{Proposition} 6.0.6.
\underline{On a l'inclusion}

\[
\alpha_S(R_S)\subset R_S(1)+pR_S.
\]

\underline{Démonstration.}
On a la factorisation (4.6.14)

\begin{equation}\tag{6.0.7}
\alpha_S=\tau^{-1}\circ\frac1p\,\psi_p\circ H(X).
\end{equation}

On sait (cf. (4.6.15))

\begin{equation}\tag{6.0.8}
H(X)\cdot R_S\subset R_S,
\end{equation}

et évidemment

\begin{equation}\tag{6.0.9}
\tau^{-1}\bigl(R_S(1)+pR_S\bigr)=R_S(1)+pR_S,
\end{equation}

de sorte qu'il suffit de vérifier

\begin{equation}\tag{6.0.10}
\frac1p\,\psi_pR_S\subset R_S(1)+pR_S,
\end{equation}

ce qui est clair par (4.6.19).

\underline{Corollaire} 6.0.11.
\underline{On a l'inclusion}

\[
\bar\alpha_S W_S\subset W_S(1).
\]

%%%%%%%%%%  scan idx 440  |  folio 433  %%%%%%%%%%

\underline{Corollaire} 6.0.12.
\underline{On a la congruence}

\[
\Delta_S(t)\equiv
\det\!\left(
1-t(\bar\alpha_S)^{\operatorname{ord}(q)}
\mathrel{\Big|}W_S(1)\right)
\pmod{\pi\mathcal O[[t]]}.
\]

\underline{Proposition} 6.1.
\underline{L'opération} $\tau^{1}$-\underline{linéaire}
$\bar\alpha_S:W_S(1)\to W_S(1)$ \underline{est égale à l'opération}

\begin{equation}\tag{6.1.1}
\tau^{-1}\psi_p\circ(X_0f(X))^{p-1}.
\end{equation}

\underline{Démonstration.}
Écrivons

\begin{equation}\tag{6.1.2}
H(X)=1+\sum_{i\geq1}H_i(X)
\end{equation}

où

\begin{equation}\tag{6.1.3}
H_i(X)\text{ est de degré $i$ en $X_0$}.
\end{equation}

D'après (4.6.19) encore, on a

\begin{equation}\tag{6.1.4}
\frac1p\,\psi_pR_S(i)\subset pR_S
\qquad\text{si }i\ne p.
\end{equation}

Compte tenu de ce que

\begin{equation}\tag{6.1.5}
H_i(X)\cdot R_S(j)\subset R_S(i+j),
\end{equation}

on a

\begin{equation}\tag{6.1.6}
\left(
\alpha_S-\frac{\tau^{-1}}p\,\psi_pH_{p-1}(X)
\right)(R_S(1))
\subset pR_S.
\end{equation}

On a donc

\begin{equation}\tag{6.1.7}
\left[
(\alpha_S)^{\operatorname{ord}(q)}
-
\left(
\frac{\tau^{-1}}p\,\psi_p\cdot H_{p-1}(X)
\right)^{\operatorname{ord}(q)}
\right]R_S(1)
\subset pR_S.
\end{equation}

Calculons maintenant $H_{p-1}(X)$. On a (cf. (4.4.4))

%%%%%%%%%%  scan idx 441  |  folio 434  %%%%%%%%%%

\begin{equation}\tag{6.1.8}
\begin{aligned}
H(X)
&=\prod_W\Theta(A_WX_0X^W)\\
&=\prod_W\exp\!\left(
\sum_{n\geq0}\frac{\pi A_WX_0X^W)^{p^n}}{p^n}
\right)\\
&=\prod_W\prod_{n\geq0}
\exp\!\left(\frac{(\pi A_WX_0X^W)^{p^n}}{p^n}\right)\\
&\equiv\prod_W\exp(\pi A_WX_0X^W)
\pmod{\text{degré }p}\\
&\equiv\exp\!\left(\pi\sum A_WX_0X^W\right)
\pmod{\text{degré }p}\\
&\equiv\exp(\pi X_0F(X))
\pmod{\text{degré }p},
\end{aligned}
\end{equation}

d'où

\begin{equation}\tag{6.1.9}
H_{p-1}(X)
\equiv\frac{\pi^{p-1}X_0^{p-1}}{(p-1)!}\,F(X)^{p-1}.
\end{equation}

On constate facilement que

\begin{equation}\tag{6.1.10}
\frac{\pi^{p-1}}p\equiv-1\pmod p
\end{equation}

donc

\begin{equation}\tag{6.1.11}
\frac1p\,H_{p-1}(X)
\equiv(X_0f(X)^{p-1}\pmod p,
\end{equation}

ce qui achève la démonstration de 6.1.

\underline{Corollaire} 6.1.12.
\underline{L'opérateur linéaire}

\[
(\bar\alpha_S)^{\operatorname{ord}(q)}:W_S(1)\longrightarrow W_S(1)
\]

\noindent
\underline{n'est autre que} $\psi_q\circ(X_0f(X))^{q-1}$.

%%%%%%%%%%  scan idx 442  |  folio 435  %%%%%%%%%%

\underline{Démonstration.}
D'après (6.1.1)

\begin{equation*}
\begin{aligned}
(\bar\alpha_S)^{\operatorname{ord}(q)}
&=
\left(
\tau^{-1}\psi_p\circ(X_0f(X))^{p-1}
\right)^{\operatorname{ord}(q)}\\
&=
\psi_q
\prod_{j=0}^{\operatorname{ord}(a)-1}
\left(X_0^{p^j}f^{\tau^j}(X^{p^j})\right)^{p-1}\\
&=
\psi_q
\prod_{j=0}^{\operatorname{ord}(q)-1}
(X_0f(X))^{p^j(p-1)}\\
&=\psi_q\circ(X_0f(X))^{q-1}.
\end{aligned}
\end{equation*}

Mettant ensemble (6.1.12) (6.0.12), et (6.0.1), on trouve,

\underline{Corollaire} 6.1.13.
\underline{On a la congruence mod} $p$ :

\begin{equation*}
Z(t,V/\mathbb F_q)
\equiv
\frac{
\det\!\left(
1-t\,\psi_q(X_0f(X))^{q-1}
\mathrel{\Big|}W_S(1)
\right)^{(-1)^{n+1}}
}{1-t}.
\end{equation*}

\subsection*{6.2.}
Il reste à calculer les groupes de cohomologie de $\mathcal O_V$.

Posons

\begin{equation}\tag{6.2.1}
\mathfrak F_q
=\text{l'opération \og élévation à la puissance $q$'ième\fg{} dans $\mathcal O_V$}.
\end{equation}

Nous avons une suite exacte de faisceaux sur
$\mathbb P^{n+1}=\mathbb P$ :

\begin{equation}\tag{6.2.2}
0\longrightarrow\mathcal O_{\mathbb P}(-d)
\xrightarrow{\,f\,}\mathcal O_{\mathbb P}
\longrightarrow\mathcal O_V\longrightarrow\mathcal O,
\end{equation}

et le diagramme suivant est commutatif

\begin{equation}\tag{6.2.3}
\begin{tikzcd}[column sep=large,row sep=large]
0 \arrow[r] &
\mathcal O_{\mathbb P}(-d) \arrow[r,"f"] \arrow[d,"f^{q-1}\mathfrak F_q"'] &
\mathcal O_{\mathbb P} \arrow[r] \arrow[d,"\mathfrak F_q"'] &
\mathcal O_V \arrow[r] \arrow[d,"\mathfrak F_q"'] &
\mathcal O\\
0 \arrow[r] &
\mathcal O_{\mathbb P}(-d) \arrow[r,"f"] &
\mathcal O_{\mathbb P} \arrow[r] &
\mathcal O_V \arrow[r] &
\mathcal O .
\end{tikzcd}
\end{equation}

%%%%%%%%%%  scan idx 443  |  folio 436  %%%%%%%%%%

Compte tenu de ce que

\begin{equation}\tag{6.2.4}
\begin{cases}
H^0(\mathbb P^{n+1},\mathcal O_{\mathbb P})\simeq\mathbb F_q\\
H^i(\mathbb P^{n+1},\mathcal O_{\mathbb P})=0&i\ne0\\
H^i(\mathbb P^{n+1},\mathcal O_{\mathbb P}(-d))=0&i\ne0,n+1,
\end{cases}
\end{equation}

la suite exacte longue de cohomologie associé à (6.2.2), donne

\begin{equation}\tag{6.2.5}
\begin{cases}
H^0(V,\mathcal O_V)\simeq\mathbb F_q\\
H^i(V,\mathcal O_V)=0&\text{si }i\ne0,n\\
H^n(V,\mathcal O_V)
\simeq H^{n+1}(\mathbb P^{n+1},\mathcal O_{\mathbb P}(-d)).
\end{cases}
\end{equation}

Le diagramme commutatif (6.2.3) donne alors

\begin{equation}\tag{6.2.6}
\det\!\left(1-t\mathfrak F_q
\mathrel{\Big|}H^0(V,\mathcal O_V)\right)=1-t
\end{equation}

et

\begin{equation}\tag{6.2.7}
\det\!\left(1-t\mathfrak F_q
\mathrel{\Big|}H^n(V,\mathcal O_V)\right)
=
\det\!\left(
1-t(f^{q-1}\mathfrak F_q)
\mathrel{\Big|}
H^{n+1}(\mathbb P^{n+1},\mathcal O_{\mathbb P}(-d))
\right..
\end{equation}

Il ne reste qu'à établir

\begin{equation}\tag{6.2.8}
\begin{aligned}
&\det\!\left(
1-t(f^{q-1}\mathfrak F_q)
\mathrel{\Big|}
H^{n+1}(\mathbb P^{n+1},\mathcal O_{\mathbb P}(-d))
\right.\\
&\qquad=
\det\!\left(
1-t\,(\psi_q\cdot(X_0f(X))^{q-1}
\mathrel{\Big|}W_S(1)
\right)
\end{aligned}
\end{equation}

ce qui est une conséquence de la

\underline{Proposition} 6.3.
\underline{Il y a un isomorphisme}

\begin{equation}\tag{6.3.1}
W_S(1)\simeq
\operatorname{Hom}_{\mathbb F_q}\!\left(
H^{n+1}(\mathbb P^{n+1}.\mathcal O_{\mathbb P}(-d),
\mathbb F_q\right)
\end{equation}

\underline{par rapport auquel l'opérateur}

\begin{equation}\tag{6.3.2}
\psi_q\circ(X_0f(X))^{q-1}:W_S(1)\longrightarrow W_S(1)
\end{equation}

%%%%%%%%%%  scan idx 444  |  folio 437  %%%%%%%%%%

\underline{est le transposé de}

\begin{equation}\tag{6.3.3}
f(X)^{q-1}\mathfrak F_q:
H^{n+1}(\mathbb P^{n+1},\mathcal O_{\mathbb P}(-d))
\longrightarrow
H^{n+1}(\mathbb P^{n+1},\mathcal O_{\mathbb P}(-d)).
\end{equation}

\underline{Démonstration.}
En calculant le groupe
$H^{n+1}(\mathbb P^{n+1},\mathcal O_{\mathbb P}(-d))$ à l'aide du
recouvrement standard de $\mathbb P^{n+1}$, on trouve

\begin{equation}\tag{6.3.4}
\begin{aligned}
H^{i+1}(\mathbb P^{n+1},\mathcal O_{\mathbb P}(-d))
&=\text{le $\mathbb F_q$-vectoriel de base des monômes}\\
&\quad X_1^{W_1}\cdots X_{n+2}^{W_{n+2}}
\qquad\text{les }W_i\in\mathbb Z\text{ vérifiant }\sum W_i=-d\\
&=\text{le $\mathbb F_q$-sous-vectoriel de base des monômes}\\
&\quad X_1^{W_1}\cdots X_{n+2}^{W_{n+2}}
\text{ qui vérifient en plus : }\exists i\text{ tel que }W_i\geq0.
\end{aligned}
\end{equation}

D'autre part, $W_S(1)=R_S(1)\otimes_{\mathcal O}\mathbb F_q$, et on a

\begin{equation}\tag{6.3.5}
\begin{aligned}
R_S(1)
&=\text{le $\mathcal O$-module libre de base les monômes}\\
&\quad
(\pi X_0)X_1^{U_1}\cdots X_{n+2}^{U_{n+2}}
\text{ vérifiant }
\sum U_i=d,\quad U_i>0\ \forall_i.
\end{aligned}
\end{equation}

L'accouplement

\begin{equation}\tag{6.3.6}
(\ ,\ ):
W_S(1)\times
H^{n+1}(\mathbb P^{n+1},\mathcal O_{\mathbb P}(-d))
\longrightarrow\mathbb F_q
\end{equation}

défini par

\begin{equation}\tag{6.3.7}
\left(
\text{classe de }\pi X_0X_1^{U_1}\cdots X_{n+2}^{U_{n+2}},
X_1^{W_1}\cdots X_{n+2}^{W_{n+2}}
\right)
=
\begin{cases}
1&\text{si }U_1=-W_i\text{ pour tout }i\\
0&\text{sinon}
\end{cases}
\end{equation}

est évidemment une dualité parfaite de $\mathbb F_q$-vectoriels, et on
constate facilement que $\psi_q\cdot(X_0f(X))^{q-1}$ est la transposée de
$f(X)^{q-1}\cdot f_q$, ce qui prouve 6.3.

%%%%%%%%%%  scan idx 445  |  folio 438  %%%%%%%%%%

\section*{Bibliographie}

\begin{enumerate}[label={[\arabic*]},leftmargin=2.4em]
\item AX, J. Zeroes of Polynomials Over Finite Fields,
  Amer. J. Math. 86, 1964, 255--261.
\item DWORK, B. On the Rationality of the Zéta Function of an
  Algebraic Variety, Amer. J. Math. 82, 1960, 631--648.
\item DWORK, B. An Algebraic Variety, in Proc. of the Woods-Hole
  Summer Institute 1964 (miméo).
\item DWORK, B. On the Zeta Function of Hypersurface II,
  Ann. Math. (2), 80, 1964, 227--299.
\item GROTHENDIECK, A., Le Groupe de Brauer III, dans
  Dix Exposés sur la Cohomologie des Schémas,
  North Holland Pub. Co, 1968.
\item SERRE, J.P. Groupes Algébriques et Corps de Classes,
  Hermann, Paris, 1959.
\item SERRE, J.P. Endomorphismes Complétement, Continus des Espaces
  de Banach $p$-Adiques, Pub. Math. I.H.E.S. n\textsuperscript{o} 12, 1962.
\item CHEVALLEY, C. Démonstration d'une hypothèse de M. Artin,
  Abhandlungen Math. Sem. Hamburg, 11, 1936, 73--75.
\end{enumerate}
